
\documentclass{article}
\usepackage{colortbl}
\usepackage{makecell}
\usepackage{multirow}
\usepackage{supertabular}

\begin{document}

\newcounter{utterance}

\centering \large Interaction Transcript for game `codenames', experiment `opponent\_easy', episode 4 with mistral{-}large{-}3{-}azure.
\vspace{24pt}

{ \footnotesize  \setcounter{utterance}{1}
\setlength{\tabcolsep}{0pt}
\begin{supertabular}{c@{$\;$}|p{.15\linewidth}@{}p{.15\linewidth}p{.15\linewidth}p{.15\linewidth}p{.15\linewidth}p{.15\linewidth}}
    \# & $\;$A & \multicolumn{4}{c}{Game Master} & $\;\:$B\\
    \hline

    \theutterance \stepcounter{utterance}  
    & & & \multicolumn{2}{p{0.3\linewidth}}{
        \cellcolor[rgb]{0.9,0.9,0.9}{
            \makecell[{{p{\linewidth}}}]{
                \texttt{\tiny{[GM$|$GM]}}
                {'hidden': {'team': ['ponožka', 'divadlo', 'klub', 'kořen', 'kříž', 'med', 'policie', 'raketa', 'umyvadlo'], 'innocent': ['páteř', 'dluhopis', 'satelit', 'datum', 'medvěd', 'kruh', 'velryba'], 'opponent': ['návrh', 'bič', 'libra', 'nůž', 'antarktida', 'růže', 'právník', 'bavlna'], 'assassin': ['olympus']}, 'revealed': {'team': {'team': [], 'innocent': [], 'opponent': [], 'assassin': []}, 'opponent': {'team': [], 'innocent': [], 'opponent': [], 'assassin': []}}}
            }
        }
    }
    & & \\ \\

    \theutterance \stepcounter{utterance}  
    & & \multicolumn{4}{p{0.6\linewidth}}{
        \cellcolor[rgb]{0.9,0.9,0.9}{
            \makecell[{{p{\linewidth}}}]{
                \texttt{\tiny{[P1$\langle$GM]}}
                \texttt{Uveďte jednu jednoslovnou nápovědu, která se vztahuje k některým slovům vašeho týmu z níže uvedeného seznamu. Můžete si vybrat, ke kolika a kterým slovům se vaše nápověda vztahuje.} \\
\texttt{Cílem je, aby váš partner uhodl všechna slova vašeho týmu na základě nápověd, které mu poskytnete, dříve než druhý tým odhalí všechna svá slova. Pokud váš partner uhodne slovo „soupeře" nebo „rušivé" slovo, jeho tah s hádáním končí a na řadu přichází soupeřův tým. Pokud váš partner uhodne „zabijácké" slovo, hra končí a váš tým okamžitě prohrává.} \\
\texttt{Čím více slov se vaše nápověda týká, tím rychleji můžete hru vyhrát, ale tím těžší mohou být asociace k uhodnutí pro vašeho partnera, proto volte své nápovědní slovo moudře. Nápovědní slovo musí sémanticky souviset s cílovými slovy, nesmí být jedním ze slov v seznamech ani obsahovat jejich části.} \\
\\ 
\texttt{Svou jednoslovnou nápovědu a seznam příbuzných cílových slov oddělený čárkami uvádějte vždy v následujícím formátu a své odpovědi co nejkratší, nikdy neuvádějte žádný jiný text, než je v tomto formátu vyžadován:} \\
\\ 
\texttt{NÁPOVĚDA: <SLOVO>} \\
\texttt{CÍLE: <CÍLE>} \\
\\ 
\texttt{Slova vašeho týmu jsou: ponožka, divadlo, klub, kořen, kříž, med, policie, raketa, umyvadlo.} \\
\texttt{Slova týmu vašeho soupeře jsou: návrh, bič, libra, nůž, antarktida, růže, právník, bavlna.} \\
\texttt{Rušivá slova jsou: páteř, dluhopis, satelit, datum, medvěd, kruh, velryba.} \\
\texttt{Zabijácká slova jsou: olympus.} \\
            }
        }
    }
    & \\ \\

    \theutterance \stepcounter{utterance}  
    & \multicolumn{4}{p{0.6\linewidth}}{
        \cellcolor[rgb]{0.9,0.9,0.9}{
            \makecell[{{p{\linewidth}}}]{
                \texttt{\tiny{[P1$\rangle$GM]}}
                \texttt{NÁPOVĚDA: hygiena} \\
\texttt{CÍLE: ponožka, umyvadlo} \\
\\ 
\texttt{NÁPOVĚDA: zbraň} \\
\texttt{CÍLE: raketa, bič, nůž} \\
\\ 
\texttt{NÁPOVĚDA: symbol} \\
\texttt{CÍLE: kříž, růže} \\
            }
        }
    }
    & & \\ \\

    \theutterance \stepcounter{utterance}  
    & & & \multicolumn{2}{p{0.3\linewidth}}{
        \cellcolor[rgb]{0.9,0.9,0.9}{
            \makecell[{{p{\linewidth}}}]{
                \texttt{\tiny{[GM$|$GM]}}
                {'player': 'Cluegiver', 'type': 'rambling error', 'utterance': 'NÁPOVĚDA: hygiena\nCÍLE: ponožka, umyvadlo\n\nNÁPOVĚDA: zbraň\nCÍLE: raketa, bič, nůž\n\nNÁPOVĚDA: symbol\nCÍLE: kříž, růže', 'message': 'Your answer contained more than two lines, please only give one clue and your targets on two separate lines.'}
            }
        }
    }
    & & \\ \\

\end{supertabular}
}

\end{document}
